% Source: https://artofproblemsolving.com/community/c3h2334551_multiplying_1_to_a_quadratic
% Stable identity: aops-topic-2334551-q01
\textbf{(a)} The reflection in the $x$-axis sends $(x,y)$ to $(x,-y)$. A point $(x,y)$ is on the graph of $f$ exactly when $y=f(x)$, i.e. when $-y=-f(x)$, i.e. when $(x,-y)$ is on the graph of $-f$. So the reflection maps one graph onto the other. The leading coefficients are $a$ and $-a$, which have opposite signs, so one parabola opens upward and the other downward.

\textbf{(b)} $f(x)=0$ holds exactly when $-f(x)=0$ (multiply or divide by $-1$). So multiplying an \emph{equation} by $-1$ does not change its solutions, even though it flips the \emph{graph} of $y=f(x)$.

\textbf{(c)} $f(x)=x^2-2x-3=(x-3)(x+1)=(x-1)^2-4$. Both parabolas cross the $x$-axis at $x=-1$ and $x=3$. The vertex of $y=f(x)$ is $(1,-4)$, and the vertex of $y=-f(x)=-(x-1)^2+4$ is its mirror image $(1,4)$.

\begin{asy}
usepackage("amsmath");
size(5cm);
defaultpen(fontsize(10pt));
import graph;
real f(real x){return x^2-2x-3;}
real g(real x){return -(x^2-2x-3);}
draw((-2.6,0)--(4.6,0),Arrow(6)); draw((0,-5.4)--(0,5.6),Arrow(6));
label("$x$",(4.6,0),E); label("$y$",(0,5.6),N);
draw(graph(f,-1.9,3.9),blue+linewidth(0.9)); draw(graph(g,-1.9,3.9),red+linewidth(0.9));
dot((-1,0)); dot((3,0)); dot((1,-4),blue); dot((1,4),red);
label("$-1$",(-1,0),NW,fontsize(8pt)); label("$3$",(3,0),NE,fontsize(8pt));
label("$(1,-4)$",(1,-4),SE,blue+fontsize(8pt)); label("$(1,4)$",(1,4),NE,red+fontsize(8pt));
label("$y=f(x)$",(3.9,f(3.9)),E,blue+fontsize(8pt)); label("$y=-f(x)$",(3.9,g(3.9)),E,red+fontsize(8pt));
\end{asy}
